% Options for packages loaded elsewhere
% Options for packages loaded elsewhere
\PassOptionsToPackage{unicode}{hyperref}
\PassOptionsToPackage{hyphens}{url}
%
\documentclass[
  11pt,
  oneside,
  10pt,
  a4paper,
  twocolumn]{article}
\usepackage{xcolor}
\usepackage[top=20mm,bottom=20mm,left=25mm,right=20mm,bottom=22mm,left=18mm,right=18mm,columnsep=6mm]{geometry}
\usepackage{amsmath,amssymb}
\setcounter{secnumdepth}{-\maxdimen} % remove section numbering
\usepackage{iftex}
\ifPDFTeX
  \usepackage[T1]{fontenc}
  \usepackage[utf8]{inputenc}
  \usepackage{textcomp} % provide euro and other symbols
\else % if luatex or xetex
  \usepackage{unicode-math} % this also loads fontspec
  \defaultfontfeatures{Scale=MatchLowercase}
  \defaultfontfeatures[\rmfamily]{Ligatures=TeX,Scale=1}
\fi
\usepackage{lmodern}
\ifPDFTeX\else
  % xetex/luatex font selection
\fi
% Use upquote if available, for straight quotes in verbatim environments
\IfFileExists{upquote.sty}{\usepackage{upquote}}{}
\IfFileExists{microtype.sty}{% use microtype if available
  \usepackage[]{microtype}
  \UseMicrotypeSet[protrusion]{basicmath} % disable protrusion for tt fonts
}{}
\makeatletter
\@ifundefined{KOMAClassName}{% if non-KOMA class
  \IfFileExists{parskip.sty}{%
    \usepackage{parskip}
  }{% else
    \setlength{\parindent}{0pt}
    \setlength{\parskip}{6pt plus 2pt minus 1pt}}
}{% if KOMA class
  \KOMAoptions{parskip=half}}
\makeatother
% Make \paragraph and \subparagraph free-standing
\makeatletter
\ifx\paragraph\undefined\else
  \let\oldparagraph\paragraph
  \renewcommand{\paragraph}{
    \@ifstar
      \xxxParagraphStar
      \xxxParagraphNoStar
  }
  \newcommand{\xxxParagraphStar}[1]{\oldparagraph*{#1}\mbox{}}
  \newcommand{\xxxParagraphNoStar}[1]{\oldparagraph{#1}\mbox{}}
\fi
\ifx\subparagraph\undefined\else
  \let\oldsubparagraph\subparagraph
  \renewcommand{\subparagraph}{
    \@ifstar
      \xxxSubParagraphStar
      \xxxSubParagraphNoStar
  }
  \newcommand{\xxxSubParagraphStar}[1]{\oldsubparagraph*{#1}\mbox{}}
  \newcommand{\xxxSubParagraphNoStar}[1]{\oldsubparagraph{#1}\mbox{}}
\fi
\makeatother


\usepackage{longtable,booktabs,array}
\newcounter{none} % for unnumbered tables
\usepackage{calc} % for calculating minipage widths
% Correct order of tables after \paragraph or \subparagraph
\usepackage{etoolbox}
\makeatletter
\patchcmd\longtable{\par}{\if@noskipsec\mbox{}\fi\par}{}{}
\makeatother
% Allow footnotes in longtable head/foot
\IfFileExists{footnotehyper.sty}{\usepackage{footnotehyper}}{\usepackage{footnote}}
\makesavenoteenv{longtable}
\usepackage{graphicx}
\makeatletter
\newsavebox\pandoc@box
\newcommand*\pandocbounded[1]{% scales image to fit in text height/width
  \sbox\pandoc@box{#1}%
  \Gscale@div\@tempa{\textheight}{\dimexpr\ht\pandoc@box+\dp\pandoc@box\relax}%
  \Gscale@div\@tempb{\linewidth}{\wd\pandoc@box}%
  \ifdim\@tempb\p@<\@tempa\p@\let\@tempa\@tempb\fi% select the smaller of both
  \ifdim\@tempa\p@<\p@\scalebox{\@tempa}{\usebox\pandoc@box}%
  \else\usebox{\pandoc@box}%
  \fi%
}
% Set default figure placement to htbp
\def\fps@figure{htbp}
\makeatother





\setlength{\emergencystretch}{3em} % prevent overfull lines

\providecommand{\tightlist}{%
  \setlength{\itemsep}{0pt}\setlength{\parskip}{0pt}}



 


\usepackage{fvextra}
\DefineVerbatimEnvironment{Highlighting}{Verbatim}{breaklines,commandchars=\\\{\}}
\usepackage{sectsty}
\allsectionsfont{\sffamily}
\usepackage{fancyhdr}
\pagestyle{fancy}
\fancyhf{}
\fancyhead[L]{Right to Read · Spend Analysis}
\fancyhead[R]{\thepage}
\renewcommand{\headrulewidth}{0.4pt}
\usepackage{xcolor}
\definecolor{red}{HTML}{DC2A14}
\definecolor{ink}{HTML}{0E0E0E}
\definecolor{cream}{HTML}{F1E8D3}
\usepackage{fvextra}
\DefineVerbatimEnvironment{Highlighting}{Verbatim}{breaklines,commandchars=\\\{\}}
% EPW-style typography: Times for body, small-caps for headings
\usepackage{mathptmx}
\usepackage[T1]{fontenc}
\usepackage{microtype}
\usepackage{fancyhdr}
\pagestyle{fancy}
\fancyhf{}
\fancyhead[L]{\small\textit{Economic \& Political Weekly} · Special Article}
\fancyhead[R]{\small\thepage}
\renewcommand{\headrulewidth}{0.4pt}
\usepackage{xcolor}
\definecolor{campaignred}{HTML}{DC2A14}
\definecolor{campaignblue}{HTML}{1936A8}
\usepackage{booktabs}
\usepackage{caption}
\captionsetup{font=small,labelfont=bf,format=plain,justification=justified}
% EPW abstract box: single-column, ruled
\usepackage{abstract}
\renewcommand{\abstractnamefont}{\normalfont\bfseries\small}
\renewcommand{\abstracttextfont}{\normalfont\small}
\setlength{\absleftindent}{0pt}
\setlength{\absrightindent}{0pt}
% Section style: small-caps, no numbering
\usepackage{sectsty}
\sectionfont{\normalsize\scshape}
\subsectionfont{\normalsize\itshape}
% Tighter paragraph spacing for EPW two-column
\setlength{\parskip}{3pt}
\setlength{\parindent}{10pt}
% Title block
\usepackage{titling}
\pretitle{\begin{flushleft}\Large\bfseries}
\posttitle{\end{flushleft}}
\preauthor{\begin{flushleft}\normalsize}
\postauthor{\end{flushleft}}
\predate{\begin{flushleft}\small\itshape}
\postdate{\end{flushleft}\vspace{-4mm}\rule{\linewidth}{0.8pt}\vspace{2mm}}
% Wide figures/tables span both columns
\usepackage{dblfloatfix}
\usepackage{stfloats}
\makeatletter
\@ifpackageloaded{caption}{}{\usepackage{caption}}
\AtBeginDocument{%
\ifdefined\contentsname
  \renewcommand*\contentsname{Table of contents}
\else
  \newcommand\contentsname{Table of contents}
\fi
\ifdefined\listfigurename
  \renewcommand*\listfigurename{List of Figures}
\else
  \newcommand\listfigurename{List of Figures}
\fi
\ifdefined\listtablename
  \renewcommand*\listtablename{List of Tables}
\else
  \newcommand\listtablename{List of Tables}
\fi
\ifdefined\figurename
  \renewcommand*\figurename{Figure}
\else
  \newcommand\figurename{Figure}
\fi
\ifdefined\tablename
  \renewcommand*\tablename{Table}
\else
  \newcommand\tablename{Table}
\fi
}
\@ifpackageloaded{float}{}{\usepackage{float}}
\floatstyle{ruled}
\@ifundefined{c@chapter}{\newfloat{codelisting}{h}{lop}}{\newfloat{codelisting}{h}{lop}[chapter]}
\floatname{codelisting}{Listing}
\newcommand*\listoflistings{\listof{codelisting}{List of Listings}}
\makeatother
\makeatletter
\makeatother
\makeatletter
\@ifpackageloaded{caption}{}{\usepackage{caption}}
\@ifpackageloaded{subcaption}{}{\usepackage{subcaption}}
\makeatother
\usepackage{bookmark}
\IfFileExists{xurl.sty}{\usepackage{xurl}}{} % add URL line breaks if available
\urlstyle{same}
\hypersetup{
  pdftitle={What India Really Spends on Public Libraries: A Revised National Expenditure Series{,} 2014--2025},
  pdfauthor={The Right to Read Campaign / CommonerLLP},
  pdfkeywords={public libraries, India, government expenditure, GDP
deflator, per-capita spending},
  hidelinks,
  pdfcreator={LaTeX via pandoc}}


\title{What India Really Spends on Public Libraries:\\
A Revised National Expenditure Series, 2014--2025}
\usepackage{etoolbox}
\makeatletter
\providecommand{\subtitle}[1]{% add subtitle to \maketitle
  \apptocmd{\@title}{\par {\large #1 \par}}{}{}
}
\makeatother
\subtitle{Real-terms per-capita analysis using official population
projections\\
and the MoSPI implicit GDP deflator}
\author{The Right to Read Campaign / CommonerLLP}
\date{2026-05-01}
\begin{document}
\maketitle
\begin{abstract}
We present a revised and extended national time series of per-capita
state expenditure on public libraries in India, covering 2014--15
through 2024--25. Our analysis builds on the expenditure compilation in
Kulkarni, Balaji \& Dhanamjaya (2025), who extracted a continuous
seven-year series from the CAG Combined Finance and Revenue Accounts
(Major Head 2205, Sub-head 105) --- the only peer-reviewed national
dataset of its kind. We introduce two methodological improvements:
first, we replace the fixed Census 2011 population denominator with
annual projected populations from the government's own Technical Group
on Population Projections (MoHFW, 2020); second, we convert nominal
expenditure to constant 2011--12 rupees using the MoSPI implicit GDP
deflator (National Accounts Statistics), the standard deflator for real
government expenditure analysis in India. We extend the series by four
years, through 2024--25, using state Demand-for-Grants documents and
CAGR extrapolation.

The revised headline is materially different from the nominal figure. At
the 2018--19 peak, India spent \textbf{₹7.87 per person per year} (real,
2011--12 ₹) on public libraries --- not ₹11.62 as the nominal,
fixed-denominator calculation suggests. By 2024--25 that figure stands
at \textbf{₹4.66}: a 41\% fall from peak in five years, and only ₹0.50
above the 2014--15 baseline of ₹4.16. The apparent nominal recovery in
the post-2020 years is largely an inflation artefact. In real terms, a
decade of budgetary activity has produced less than fifty paise per
person per year of net progress.
\end{abstract}


\subsection{Introduction}\label{introduction}

India does not publish consolidated data on public library spending.
State-level allocations are embedded, without headline visibility, in
the CAG's Combined Finance and Revenue Accounts under Major Head 2205
(Education, Sports, Art and Culture), Sub-head 105 (Public Libraries).
No ministry reports an annual aggregate. No planning document tracks the
per-capita trend. In this context, Kulkarni, Balaji \& Dhanamjaya (2025)
performed a foundational service: their systematic extraction of MH
2205-105 across states and years produced what is, to our knowledge, the
first peer-reviewed continuous national time series for this budget
head.

The policy implications of any such series depend on whether the numbers
accurately capture the real-terms trajectory of public investment. Two
features of the original analysis limit that accuracy. First, per-capita
figures are computed against India's Census 2011 total population, held
constant across all years. India's population has grown by approximately
170 million since that enumeration; the fixed denominator progressively
flatters the per-capita figure as years advance. Second, the series is
expressed in nominal rupees --- appropriate as a record of what
legislatures voted, but not as a measure of what those votes actually
purchased, or of how spending power has changed over time.

These are not incidental concerns. The decade covered (2014--15 to
2024--25) has seen India's GDP deflator fall from 0.84 to 0.57 relative
to a 2011--12 base --- a 32\% compression in the real value of one
current-price rupee. Any analysis that ignores this will, over a
ten-year window, produce apparent growth that is partly or largely
illusory.

This paper introduces the two corrections, presents a revised national
series, extends it through 2024--25, and discusses the changed policy
picture that results.

\subsection{Data and Methods}\label{data-and-methods}

\subsubsection{Expenditure base}\label{expenditure-base}

The state-level expenditure figures for 2014--15 through 2020--21 are
taken directly from Kulkarni, Balaji \& Dhanamjaya (2025), Table 1: CAG
Combined Finance and Revenue Accounts, MH 2205-105, state-level actuals.
We build on their compilation and do not re-extract from primary CAG
sources for this period.

For 2021--22 through 2024--25, CAG actuals are unavailable or not yet
published. We extend using two methods in order of preference: (1) state
Demand-for-Grants documents where available --- Assam, Goa, Rajasthan
and Odisha provide usable figures for selected post-2021 years; and (2)
CAGR extrapolation for remaining states, applying the compound annual
growth rate computed over 2016--17 through 2020--21 year-on-year from
the 2020--21 base. This window excludes the 2015--16 base-year effect
and captures the pre-COVID trend.

States whose series show anomalous year-on-year volatility attributable
to cess-reclassification artefacts (West Bengal, Karnataka, Telangana,
Tripura) receive CAGR treatment with a flagged caveat.

\subsubsection{Population denominator}\label{population-denominator}

The original analysis holds India's Census 2011 total population
constant across all years. We substitute the \textbf{Technical Group on
Population Projections for India and States 2011--2036} (National
Commission on Population, MoHFW, July 2020), Table 11 --- annual July-1
projected totals. These are the official government projections used in
Fifteenth Finance Commission resource-devolution calculations and major
health programme planning.

By 2018--19, the Census 2011 fixed denominator understates the projected
national population by approximately 10\%. By 2024--25 the gap grows to
approximately 15\%. Using the fixed denominator assigns accumulated
population growth as apparent per-capita gain, introducing a systematic
upward bias that widens over the series.

\subsubsection{Price deflator}\label{price-deflator}

We deflate using the \textbf{MoSPI implicit GDP deflator}, derived from
the National Accounts Statistics GDP-at-constant-prices and
GDP-at-current-prices series (base year 2011--12). This is the standard
for real government expenditure analysis in India, used by the Reserve
Bank of India's annual State Finances study, the OECD Government at a
Glance real expenditure tables, IMF Article IV reviews, Finance
Commission reports (most recently the Fifteenth), and CAG performance
audits. The Consumer Price Index for Industrial Workers (CPI-IW),
designed for dearness-allowance calculations for urban industrial
workers, is not the appropriate deflator for government-spending
analysis. We do not use it.

The deflator factors applied to each year in the series are presented in
Table~\ref{tbl-deflators}.

\begin{table}

\caption{\label{tbl-deflators}MoSPI implicit GDP deflator factors
(constant 2011--12 ÷ current prices). Source: National Accounts
Statistics, MoSPI.}

\centering{

\begin{verbatim}
| Year    |   Deflator factor | NAS revision stage   |
|:--------|------------------:|:---------------------|
| 2014-15 |            0.8444 | Third Revised        |
| 2015-16 |            0.8256 | Third Revised        |
| 2016-17 |            0.7997 | Third Revised        |
| 2017-18 |            0.7691 | Third Revised        |
| 2018-19 |            0.7404 | Third Revised        |
| 2019-20 |            0.723  | Third Revised        |
| 2020-21 |            0.6898 | Third Revised        |
| 2021-22 |            0.6366 | Second Revised       |
| 2022-23 |            0.6011 | Final Estimates      |
| 2023-24 |            0.586  | First Revised        |
| 2024-25 |            0.5684 | Provisional          |
\end{verbatim}

}

\end{table}%

National-level deflator factors are applied uniformly across all states,
consistent with Finance Commission practice.

\subsubsection{The three series}\label{the-three-series}

We construct three national per-capita series for direct comparison
throughout this paper:

\begin{itemize}
\tightlist
\item
  \textbf{Series A} (Balaji et al.~method): Nominal total expenditure ÷
  Census 2011 fixed population. Replicates the original analysis.
\item
  \textbf{Series B} (population correction only): Nominal total
  expenditure ÷ TG 2020 annual projected population. Isolates the
  denominator effect.
\item
  \textbf{Series C} (both corrections --- our preferred series):
  Expenditure deflated to constant 2011--12 ₹ via the NAS implicit GDP
  deflator ÷ TG 2020 annual projected population.
\end{itemize}

Series C is the measure we recommend for per-capita real public
investment analysis.

\subsection{Results}\label{results}

\subsubsection{National trend: three series
compared}\label{national-trend-three-series-compared}

Table~\ref{tbl-series} presents the full numerical results.
\textbf{?@fig-trend} plots the three series together, with the
solid/dashed distinction marking the boundary between CAG actuals
(through 2020--21) and projections (2021--22 onwards).

\begin{table}

\caption{\label{tbl-series}National per-capita public library
expenditure (₹ per person per year), three methods, 2014--15 to
2024--25. † Projection.}

\centering{

\begin{verbatim}
| Year     |   A: Nominal, Census 2011 |   B: Nominal, TG pop |   C: Real 2011–12 ₹, TG pop |
|:---------|--------------------------:|---------------------:|----------------------------:|
| 2014-15  |                      5.16 |                 4.93 |                        4.16 |
| 2015-16  |                      6.62 |                 6.25 |                        5.16 |
| 2016-17  |                      8.62 |                 8.05 |                        6.44 |
| 2017-18  |                      9.96 |                 9.19 |                        7.07 |
| 2018-19  |                     11.64 |                10.63 |                        7.87 |
| 2019-20  |                      7.79 |                 7.04 |                        5.09 |
| 2020-21  |                      8.39 |                 7.5  |                        5.17 |
| 2021-22† |                      8.49 |                 7.51 |                        4.78 |
| 2022-23† |                      8.66 |                 7.6  |                        4.57 |
| 2023-24† |                      9.13 |                 7.94 |                        4.65 |
| 2024-25† |                      9.51 |                 8.2  |                        4.66 |
\end{verbatim}

}

\end{table}%

Three findings stand out from \textbf{?@fig-trend} and
Table~\ref{tbl-series}.

\textbf{Population correction (A → B).} Substituting TG 2020 projections
for the Census 2011 fixed denominator reduces the 2018--19 peak from
₹11.64 to ₹10.63 --- a 9\% revision. By 2024--25, Series A shows ₹9.51
versus Series B's ₹8.20. This gap will continue widening until a new
census resets the denominator.

\textbf{Real-terms correction (B → C).} This is the larger revision.
Deflating Series B's 2018--19 peak of ₹10.63 to constant 2011--12 ₹
yields ₹7.87 --- a further 26\% reduction. The nominal upward trend
visible after 2020--21 in Series A and B is absorbed almost entirely by
inflation: Series C moves from ₹5.17 (2020--21) to ₹4.66 (2024--25), a
near-flat line.

\textbf{The extended series.} After 2020--21, real per-capita library
expenditure (Series C) traces ₹4.78, ₹4.57, ₹4.65, ₹4.66. There is no
recovery. The 2024--25 real figure is only ₹0.50 above the 2014--15
starting point of ₹4.16. Ten years of nominal budget growth has produced
less than fifty paise of real-terms progress per person per year.

\subsubsection{The two corrections: a
decomposition}\label{the-two-corrections-a-decomposition}

\textbf{?@fig-decomp} illustrates the size of each correction
independently over the series. The gap between Series A and Series C can
be split into (i) the denominator effect and (ii) the deflation effect.

The decomposition shows that the inflation effect dominates throughout
the series and grows over time, while the population effect is secondary
and roughly constant in absolute terms. By 2018--19, the total upward
bias in the original series is ₹3.77 per person (₹11.64 − ₹7.87); of
this, ₹2.76 is attributable to inflation and ₹1.01 to the fixed
denominator.

\subsubsection{State cross-section (2018--19, real
terms)}\label{state-cross-section-201819-real-terms}

\textbf{?@fig-state} converts the 2018--19 state-level nominal figures
from Balaji et al.~to real 2011--12 ₹ using the national deflator factor
of 0.7404. All 31 states and UTs are ranked.

The cross-section reveals a distribution that nominal figures soften.
The highest-spending state (Goa) spends ₹87.93 per person per year in
real terms; the lowest (Jharkhand) spends ₹0.10. The gap is 879-fold.
The median state spends approximately ₹8 per person per year. The four
Central Zonal Council states --- Uttar Pradesh, Madhya Pradesh,
Chhattisgarh, Uttarakhand --- fall below ₹2 per person in real terms;
together they account for over a quarter of India's population.

\subsubsection{The cess paradox: real versus nominal
volatility}\label{the-cess-paradox-real-versus-nominal-volatility}

Five states incorporate a statutory library cess --- Tamil Nadu, Andhra
Pradesh, Karnataka, West Bengal, and Kerala. Their nominal series show
substantially more year-on-year volatility than their real series
(\textbf{?@fig-cess}), suggesting that a portion of the apparent nominal
variation reflects price-level changes rather than genuine budget
decisions.

Karnataka's apparent nominal spike in 2017--18 --- from approximately
₹9,800 lakh to ₹23,100 lakh in a single year --- deflates to ₹18,500
lakh in real terms: still significant, but more legible as an actual
policy event once the price-level effect is removed.

\subsection{Discussion}\label{discussion}

\subsubsection{The nominal narrative versus the real
one}\label{the-nominal-narrative-versus-the-real-one}

The Balaji et al.~series, as published, supports a narrative of
substantial growth followed by disruption: from ₹5.16 in 2014--15 to a
peak of ₹11.64 in 2018--19, a fall to ₹7.79 in 2019--20, partial
recovery by 2020--21, and continued recovery implied by the nominal
trend thereafter. The shape of the story is growth, shock, recovery.

The corrected Series C tells a different story. Real per-capita spending
rose from ₹4.16 to ₹7.87 between 2014--15 and 2018--19 --- a genuine
gain, but one driven by nominal budget expansion into a period of
relatively slower aggregate price growth, not by structural policy
commitment. Then a sharp fall to ₹5.09 in 2019--20, which preceded any
COVID impact on state budgets. From 2021--22, a slow slide to ₹4.66 ---
not recovery, but consolidation near the 2014--15 baseline. The shape of
the real story is a shallow gain, a structural withdrawal, a decade
closed near where it opened.

The policy implication differs fundamentally. The nominal narrative
suggests the task is to recover the 2018--19 level. The real-terms
narrative establishes that the 2018--19 level was itself insufficient
--- ₹7.87 in constant 2011--12 terms, against a constitutional mandate
for universal public library access --- and that the structural
direction since 2018--19 has been downward regardless of nominal budget
headlines.

\subsubsection{Why the GDP deflator and not
CPI-IW}\label{why-the-gdp-deflator-and-not-cpi-iw}

The choice of deflator matters numerically. CPI-IW (Consumer Price Index
for Industrial Workers, Labour Bureau) is a narrower basket and moves
faster than the implicit GDP deflator. Using CPI-IW would produce lower
real estimates across the board. We chose the GDP deflator because it
reflects the price level of all goods and services produced in the
economy, including government services. Government library expenditure
purchases staff time, books, building maintenance, and equipment --- a
basket better approximated by the general economy's price level than by
a basket designed for urban industrial workers' consumption. This is the
explicit rationale in the IMF Government Finance Statistics Manual
(2014), and it is the deflator used by the RBI in its annual State
Finances study. We follow that standard.

\subsubsection{Why annual projected
denominators}\label{why-annual-projected-denominators}

Using a fixed census denominator is common in Indian administrative data
analysis because census enumeration is the only comprehensive population
count available and the 2021 Census has not yet been published. But the
right response to that gap is to use the government's own official
projections --- the TG 2020 series --- produced precisely for this
purpose and already used by the Finance Commission and major health
programme evaluations.

There is a further reason to prefer annual denominators: they prevent
the fixed-denominator artefact, by which India's 1.2\% annual population
growth registers as apparent per-capita spending growth even when
aggregate expenditure is flat. Eliminating this artefact is a
precondition for an honest multi-year comparison.

\subsection{Conclusion}\label{conclusion}

This paper presents a revised and extended national per-capita series
for Indian public library expenditure, building on the dataset of
Kulkarni, Balaji \& Dhanamjaya (2025). Two methodological improvements
--- annual population projections replacing a fixed census denominator,
and real-terms deflation using the MoSPI implicit GDP deflator ---
change both the headline figures and the policy conclusion. Both
corrections use the Indian government's own published data.

The corrected figures for the period covered by the original analysis
are:

{\def\LTcaptype{none} % do not increment counter
\begin{longtable}[]{@{}
  >{\raggedright\arraybackslash}p{(\linewidth - 4\tabcolsep) * \real{0.2727}}
  >{\raggedleft\arraybackslash}p{(\linewidth - 4\tabcolsep) * \real{0.3636}}
  >{\raggedleft\arraybackslash}p{(\linewidth - 4\tabcolsep) * \real{0.3636}}@{}}
\toprule\noalign{}
\begin{minipage}[b]{\linewidth}\raggedright
Year
\end{minipage} & \begin{minipage}[b]{\linewidth}\raggedleft
Original (nominal, Census 2011)
\end{minipage} & \begin{minipage}[b]{\linewidth}\raggedleft
This paper (real 2011--12 ₹, TG 2020)
\end{minipage} \\
\midrule\noalign{}
\endhead
\bottomrule\noalign{}
\endlastfoot
2014--15 & ₹5.16 & \textbf{₹4.16} \\
2018--19 (peak) & ₹11.64 & \textbf{₹7.87} \\
2020--21 & ₹8.39 & \textbf{₹5.17} \\
\end{longtable}
}

Extended through 2024--25, the real-terms per-capita figure stands at
\textbf{₹4.66} --- 41\% below peak, and only marginally above the
decade's starting point. A nominal upward drift in the post-2020 years
is absorbed by inflation, leaving no recovery in real terms.

The policy picture that emerges from the corrected series is not one of
temporary disruption from which recovery is underway. It is one of a
shallow real-terms gain from 2014 to 2018, a structural withdrawal that
preceded the pandemic, and a decade that closes near where it opened.
Nominal budget increases over this period have been substantially
consumed by inflation and population growth, leaving a per-person real
investment in public libraries that is lower today than at any point
since 2016--17.

We recommend that future analyses of Indian public library expenditure
adopt (1) TG 2020 annual projected population as denominator until the
2021 Census is released, and (2) the MoSPI implicit GDP deflator as the
deflation standard, consistent with Finance Commission, RBI, and IMF
practice.

\subsection{References}\label{references}

Kulkarni, S. R., Balaji, B. P., \& Dhanamjaya, M. (2025). {[}Full
citation to be inserted --- verify journal and title before
submission.{]}

Ministry of Health and Family Welfare, Government of India. (2020).
\emph{Population Projections for India and States 2011--2036.} Technical
Group on Population Projections, National Commission on Population. New
Delhi: MoHFW. URL:
\texttt{nhm.gov.in/New\_Updates\_2018/Report\_Population\_Projection\_2019.pdf}.
{[}Table 11.{]}

Ministry of Statistics and Programme Implementation, Government of
India. (annual). \emph{National Accounts Statistics.} New Delhi: MoSPI.
{[}GDP at current and constant 2011--12 prices; Third Revised through
Provisional estimates.{]}

Reserve Bank of India. (2026). \emph{State Finances: A Study of Budgets
of 2025--26.} Mumbai: RBI. {[}Statement 31; real-terms methodology.{]}

International Monetary Fund. (2014). \emph{Government Finance Statistics
Manual 2014.} Washington, D.C.: IMF. {[}Chapter 10: deflation standards
for government expenditure series.{]}

\begin{center}\rule{0.5\linewidth}{0.5pt}\end{center}

\emph{Derivation scripts and input data: \texttt{spend/spend\_data.py}
and \texttt{spend/index.qmd} in the project repository. All figures are
reproducible from the cited sources.}




\end{document}
